\begin{figure}[t]
\centering
    \begin{subfigure}{.49\textwidth}
      \centering
      \includegraphics[width=\linewidth]{media/stage3/terrain-walker-images/left_wall-square.png}
      \caption{\textsc{left-wall}}
      \label{fig:left-wall}
    \end{subfigure}
    \hfil
    \begin{subfigure}{.49\textwidth}
      \centering
      \includegraphics[width=\linewidth]{media/stage3/terrain-walker-images/right_wall-monster.png}
      \caption{\textsc{right-wall}}
      \label{fig:right-wall}
    \end{subfigure}
    \\
    \begin{subfigure}{.49\textwidth}
      \centering
      \includegraphics[width=\linewidth]{media/stage3/terrain-walker-images/bumpy-wheel.png}
      \caption{\textsc{bumpy}}
      \label{fig:bumpy}
    \end{subfigure}
    \hfil
    \begin{subfigure}{.49\textwidth}
      \centering
      \includegraphics[width=\linewidth]{media/stage3/terrain-walker-images/tunnel1.png}
      \caption{\textsc{tunnel}}
      \label{fig:tunnel}
    \end{subfigure}%
    \caption{\textbf{Terrains used in experiments}. A small set of terrains from which distributions that force \emph{conditionality} can be constructed. The terrains are (a) \textsc{left-wall}, (b) \textsc{right-wall}, (c) \textsc{bumpy}, and (d) \textsc{tunnel}. The Sodaracers produced by the models are incapable of scaling the walls in \textsc{left-wall} and \textsc{right-wall}, and therefore must produce different Sodaracers for these two terrains. Similarly, achieving good performance in the \textsc{tunnel} terrain can only be achieved with short Sodaracers, which struggle to locomote as quickly as taller ones, encouraging the model to distinguish between these terrains. Finally, Sodaracers that are proficient in locomotion on flat terrains tend to perform poorly on \textsc{bumpy}, encouraging the model to produce yet another Sodaracer for this terrain. In contrast to \textsc{tunnel}, which requires Sodaracers with a particular morphology, achieving good performance on \textsc{bumpy} requires modifying the way Sodaracers \emph{locomote}. Example Sodaracers are added to the figures to illustrate the scale of the terrains.}
    \label{fig:simple-terrain-distribution}
\end{figure}
